
\DeclareMathOperator{\Cov}{Cov}
\DeclareMathOperator{\Var}{Var}
\DeclareMathOperator{\E}{\mathbb{E}}
\DeclareMathOperator{\Proba}{\mathbb{P}}

\newcommand{\Covb}[2]{\ensuremath{\Cov\!\left[#1,#2\right]}}
\newcommand{\Eb}[1]{\ensuremath{\E\!\left[#1\right]}}
\newcommand{\Pb}[1]{\ensuremath{\Proba\!\left[#1\right]}}
\newcommand{\Varb}[1]{\ensuremath{\Var\!\left[#1\right]}}

% norm
%\newcommand{\norm}[1]{\| #1 \|}

\newcommand{\indep}{\rotatebox[origin=c]{90}{$\models$}}





\usepackage{mathptmx,amsmath,amssymb,graphicx,bibentry,bbm,babel,ragged2e}

\makeatletter


\renewcommand{\addlogo}{
	\hfill	
	\includegraphics[height=.4cm]{figures/logos/iconOM.png}
	\includegraphics[scale=1]{figures/logos/openmole_dark.png}
}


\newcommand{\noun}[1]{\textsc{#1}}
\newcommand{\jitem}[1]{\item \begin{justify} #1 \end{justify} \vfill{}}
%\newcommand{\sframe}[2]{\frame{\frametitle{#1}\addlogo #2}}

\newenvironment{centercolumns}{\begin{columns}[c]}{\end{columns}}
%\newenvironment{jitem}{\begin{justify}\begin{itemize}}{\end{itemize}\end{justify}}

%\usetheme{Warsaw}
%\setbeamertemplate{footline}[text line]{}
%\setbeamersize{text margin left=15pt,text margin right=15pt}
%\setbeamertemplate{headline}{}
%\setbeamertemplate{footline}[frame number]
%\setbeamertemplate{navigation symbols}{}

\usetheme{Darmstadt}
\setbeamertemplate{headline}{}
\setbeamertemplate{navigation symbols}{}
\setbeamercolor{palette quaternary}{fg=primaryDarkOM, bg=primaryGreenOM}


%\setbeamercovered{transparent}
%\setbeamercolor{structure}{fg=purple!50!blue, bg=purple!50!blue}

\@ifundefined{showcaptionsetup}{}{%
 \PassOptionsToPackage{caption=false}{subfig}}
\usepackage{subfig}

\usepackage[utf8]{inputenc}
%\usepackage[T1]{fontenc}


\usepackage{tikz}

\usepackage{multirow}

\usepackage{mdframed}

%\usepackage[usenames,dvipsnames]{pstricks}
%\usepackage{auto-pst-pdf}


%\usepackage[dvipsnames]{xcolor}

\usepackage{threeparttable}


\usepackage{listings}
\lstset{language=Java} 

\makeatother




\title{A roadmap for spatial sensitivity analysis}
\author{}
\date{}


\maketitle

\justify


\begin{abstract}

\end{abstract}





%%%%%%%%%%%%
\section{Introduction}


\subsection{Sensitivity analysis methods}

\cite{lilburne2009sensitivity} study how standard sensitivity analysis methods can be applied to spatial models, and show that Sobol sensitivity index is the most suited for such models. \cite{saint2010comparison} benchmark three sensitivity analysis techniques and also concludes that Sobol is better in reflecting changes in model behavior.

In environmental science, multi-criteria decision analysis based on GIS has been tested for the sensitivity to weight \cite{chen2010spatial}.



\subsection{Synthetic data}

\cite{cottineau2017initial} suggest the use of synthetic data generation to test the sensitivity of simple agent-based models to their spatial initial configuration.

\cite{obled1994sensitivity} combine observational data with a baseline model to understand to sensitivity of hydrographs to the spatial distribution of rainfall.

Link with methods in geostatistics \cite{gotway2002combining} (extrapolation generates synthetic data ?)


\subsection{Real data}

\cite{Thomasetal2017} show that city boundary selection can have a significant impact on model behavior for Luti models.

\cite{moderl2011spatial} sensitivity of infrastructure network to spatial distribution of hazards






%%%%%%%%%%%%%%
\subsection{Literature mapping}

% Initial corpus : keyword request + screening
% "Spatial sensitivity analysis"
% "Spatial synthetic data" -> bof
% "Spatial noise" -> bof
%





%%%%%%%%%%%%
\section{Methods}


\subsection{General view}

Two global axis can be formulated:

\begin{enumerate}
	\item The generation of synthetic spatial data
	\item The import and manipulation of real spatial data 
\end{enumerate}

A longer term goal is a stronger link between these two. The different axis belong to different knowledge domains \cite{raimbault2017applied} and their relations foster integrative approaches.



\subsection{Generation of synthetic spatial data}


\subsubsection{Raster data}

\begin{itemize}
	\item Single layer synthetic grids ; several methods implemented :
		\begin{itemize}
			\item Kernel mixtures \cite{anas1998urban}
			\item Reaction-diffusion model generating population distributions \cite{raimbault2018calibration}
		\end{itemize}
	\item Multi-layer synthetic grids
	\item Multi-layer with specified correlation structure
\end{itemize}


\paragraph{Kernel mixtures}


\paragraph{Reaction-diffusion}





\subsubsection{Vector data}

\begin{itemize}
	\item Transportation networks
	\item Spatialized `classic` network generation models
	\item Correlated layers	
	\item Correlated network and raster 
\end{itemize}



\subsection{Perturbation of real data}

\cite{constantine2012method} example of Gaussian field perturbation







%%%%%%%%%%%%%%%%%%
\section{Implementation}


\subsection{Architecture}

\begin{itemize}
	\item Methods available as a standalone library
	\item Integrated into OpenMOLE as a plugin
\end{itemize}



\subsection{Example of a workflow}

We illustrate spatial sensitivity analysis by coupling a spatial generator embedded into OpenMOLE as a \texttt{Sampling} with a toy Luti model, using new features of the \texttt{NetLogoTask}

\begin{itemize}
	\item Generate realistic correlated spatial grids with an exponential mixture using a \texttt{SpatialSampling}
	\item Plug this into a NetLogoTask (recursive arrays inputs) for a toy Luti model
	\item Get successive system states (recursive array outputs)
\end{itemize}




%%%%%%%%%%%%%%%%%%%%
\section{Application}

First application : comparison of the scala implementation and the netlogo implementation of the Luti part of the Lutecia model (only spatial setup is random).







%%%%%%%%%%%%%%%%%%%%
%% Biblio
%%%%%%%%%%%%%%%%%%%%

\bibliographystyle{apalike}
\bibliography{biblio.bib}


\end{document}
